\section{Related Work}
\label{sec:related}
\paragraph{Causal skill attribution.}
Paired trace auditing, matched/no-skill comparisons, and randomized skill masking already show that exposing an agent to a skill can causally change behavior and that per-skill effects are heterogeneous~\citep{zhou2026counterfactual,dong2026agentskills,wang2026notallskills}. We therefore do not claim paired skill causality, no-skill attribution, or heterogeneous skill effects as new. Our residual question is stricter: for a \emph{predeclared recurring failure}, does one targeted operation repair the original agent while also outperforming matched generic assistance? The observed T$>$G but T$=$N cutoff cell shows why the second comparison changes the scientific verdict rather than merely adding another baseline.

\paragraph{Skill evolution and reusable procedures.}
ContinualSkillBench, Evo-Harness, SkillProx, HyperSkill, and Anything2Skill study continual skill maintenance, compilation, optimization, structured skill memory, and conversion of external knowledge into reusable procedures~\citep{guan2026continualskillbench,wei2026evoharness,zheng2026skillprox,xu2026hyperskill,pan2026anything2skill}. These works make it untenable to claim that procedure reuse, cross-task skill accumulation, or external-knowledge-to-skill compilation is itself novel. We instead treat skill acquisition as upstream of our object: once a reusable procedure exists, our three-arm contract asks whether it is a \emph{binding repair} for the recurring failure that motivated its persistence, and unchanged-code replay tests where that diagnosis transfers or disappears at ceiling.

\paragraph{Temporal retrieval and generic assistance.}
Selective/corrective RAG and temporal retrieval systems already address irrelevant evidence, time-sensitive indexing, and explicit temporal constraints upstream of generation~\citep{asai2024selfrag,yan2024crag,zhang2025mrag,abdallah2026tempretriever,han2025temporalgraphs,wang2026iarag,liu2026chronos}. A useful cutoff or alignment operation therefore need not be packaged as a callable skill. We test this directly rather than leaving it as a limitation. Our R treatment runs the same frozen T operation on the same candidate evidence and places the identical output in retrieved context, while T exposes it through callable state. On the 18 pre-frozen DeepSeek load-bearing endpoints T--R is 0.0 points with a 90\% interval of $[-5.6,5.6]$, and the secondary Kimi alignment check is also exactly tied. We therefore explicitly concede that temporal retrieval/context construction is an equally effective integration surface at the tested resolution. Likewise, G$_0$ shows that the target-blind G deficits are not an unavoidable cost of adding a helper. The residual scientific object is consequently \emph{procedure-bottleneck identification}: which temporal operation is binding in which model--task regime, not a new temporal-retrieval algorithm and not evidence that skill packaging itself is necessary.
